\documentclass[10pt,a4paper]{article}
\usepackage{amsmath,amssymb,amsthm,graphicx}
% Card-local minimal preamble (journal unit sigma-2023-019-higher-braidings-of-diagonal-type, P2-ahl16 precedent).
%% Verbatim from the pinned source: the author theorem environments and macro block,
%% so every macro and environment of the counted spans is available.
% Omitted and why: the documentclass line (the journal class sigma is not on the host;
%% the assembler supplies the class) and page-geometry/font blocks (they fight the
%% forced xelatex engine).
\usepackage{amsmath,amssymb,amsthm,mathtools}

\numberwithin{equation}{section}
\newtheorem{Theorem}{Theorem}[section]
\newtheorem*{Theorem*}{Theorem}
\newtheorem{Corollary}[Theorem]{Corollary}
\newtheorem{Lemma}[Theorem]{Lemma}
\newtheorem{Proposition}[Theorem]{Proposition}
\newtheorem{conje}[Theorem]{Conjecture}
 { \theoremstyle{definition}
\newtheorem{Definition}[Theorem]{Definition}
\newtheorem{Note}[Theorem]{Note}
\newtheorem{Example}[Theorem]{Example}
\newtheorem{Remark}[Theorem]{Remark}
\newtheorem{conve}[Theorem]{Convention}
}
\newcommand{\NN }{\mathbb{N}}
\newcommand{\CC }{\mathbb{C}}
\newcommand{\RR }{\mathbb{R}}
\newcommand{\FF }{\mathbb{F}}
\newcommand{\QQ }{\mathbb{Q}}
\newcommand{\ZZ }{\mathbb{Z}}
\newcommand{\id }{\mathrm{id}}
\DeclareMathOperator{\Aut}{Aut}
\DeclareMathOperator{\End}{End}
\DeclareMathOperator{\Hom}{Hom}
\DeclareMathOperator{\GL}{GL}
\DeclareMathOperator{\SL}{SL}
\DeclareMathOperator{\rank}{rank}
\newcommand{\Cc }{\mathcal{C}}
\newcommand{\Wg }{\mathcal{W}}
\newcommand{\rsC }{\mathcal{R}}
\newcommand{\re }{^\mathrm{re}}
\newcommand{\rer }[1]{(R\re)^{#1}}
\DeclareMathOperator{\Mor}{Mor}
\DeclareMathOperator{\Obj}{Obj}
\DeclareMathOperator{\md}{mod}
\newcommand{\qq }{{\bf q}}
\newcommand{\KK }{\mathbb{K}}
\DeclareMathOperator{\Irr}{Irr}
\newcommand{\del }{\partial^L}
\newcommand{\der }{\partial^R}
\newcommand{\BB }{\mathfrak{B}}
\newcommand{\II }{\mathfrak{I}}
\newcommand{\Kc }{\mathcal{K}}
\DeclareMathOperator{\ad}{ad}
\newcommand{\dn }{d}
\newcommand{\qR }{\tilde R}
\newcommand{\bich }{\chi}
\usepackage{cleveref}
\crefname{originalproposition}{original proposition}{original propositions}
\Crefname{originalproposition}{Original proposition}{Original propositions}
\makeatletter
\newlabel{RdCartan}{{3.8}{source}{}{sourcepointer-1}{}}
\newlabel{RdCartan@cref}{{[originalproposition][0][]3.8}{[1][0][]source}{}{}{}}
\makeatother
\crefname{originaltheorem}{original theorem}{original theorems}
\Crefname{originaltheorem}{Original theorem}{Original theorems}
\makeatletter
\newlabel{cartan_d}{{3.7}{source}{}{sourcepointer-2}{}}
\newlabel{cartan_d@cref}{{[originaltheorem][0][]3.7}{[1][0][]source}{}{}{}}
\makeatother
\title{Closure under the tensor-power reflections makes a tensor set into a Cartan graph}
\author{Cuntz, Michael, Ohrmann, Tobias}
\date{doi:10.3842/SIGMA.2023.019; CC BY 4.0}
\begin{document}
\maketitle
\setcounter{section}{3}
\setcounter{Theorem}{8}
\setcounter{equation}{8}
\noindent\textit{Editorial note.} All delivered mathematical text is the author's own, copied from the pinned journal source file sources/sigma-2023-019-higher-braidings-of-diagonal-type/pinned-original-source.tex (Higher Braidings of Diagonal Type); no statement or proof step was rewritten, reordered or removed.\par\medskip
\noindent\textit{Editorial note.} The counted proof is the author's complete argument: the author verifies the two Cartan-graph axioms for the reflected set A: (C1) from Proposition RdCartan plus the order-two property of the defining reflection, (C2) by comparing Cartan entries through the R\_m factorization and the eigenvector relations of (sigma\textasciicircum{}q\_l)\textasciicircum{}\{tensor dn\}, closing via eqref cartan\_d. In-span labels: the proposition itself renders as Proposition 3.9; the span carries no numbered display environments. Out-of-span navigation pointers declared as external references: RdCartan (Theorem 3.8, line 734, used at line 756) and cartan\_d (Theorem 3.7, line 714, used at line 768 via \textbackslash{}eqref). The factorization formula quoted at line 761 is the content of Rossod (3.8), restated inline inside the proof. Seed convention: counters stepped from the source itself and cross-validated against a full host compile of the pinned source (aux ground truth): section 3 of the shared [Theorem] counter holds seven numbered items before this proposition (3.1 through 3.8 with RdCartan = 3.8 and cartan\_d = 3.7), and section 3 holds numbered equations (3.1) through (3.8) before the span, so the card prints Proposition 3.9.\par\medskip
\noindent\textit{Editorial note.} The label RdGraph names the counted statement, and the proof environment immediately following it in the source is delivered verbatim as its proof.\par\medskip
\noindent\textit{Editorial note.} The counted statement is carried by the author's own Proposition environment in the article source.\par\medskip
\noindent\textit{Editorial note.} Journal-unit assembly: card-local preamble (author theorem environments and macros verbatim), article class, no bibliography in the counted span (the span carries no \textbackslash{}cite commands).\par\medskip
\noindent\textit{Editorial scope.} the counted statement and its own proof environment only\par\medskip
\noindent\textit{Allowed standard prerequisites (editorial).}\begin{itemize}
\item Cartan graphs: a quadruple (I, A, (rho\_i), (c\textasciicircum{}q)) with finite index set I, set of objects A, reflections rho\_i, and generalized Cartan matrices c\textasciicircum{}q, satisfying axioms (C1) and (C2) (source Section 3 definitions).
\item The tensor-power construction of Section 3: for a tensor q in K\textasciicircum{}\{I\textasciicircum{}dn\} with defined Cartan matrix c\textasciicircum{}q, the reflection operator sigma\textasciicircum{}q\_l and its eigenvectors v\_m, w\_m, s\_m with the eigenvalue relations (sigma\textasciicircum{}q\_l)\textasciicircum{}\{tensor dn\}(v\_m) = v\_m, (w\_m) -> -w\_m, (s\_m) -> s\_m (source Section 3.2, recalled inside the proof at lines 762-764).
\item The factorization Rossod (3.8) of the condition R\_m = 0 as (1 - bich\_q(v\_m))(1 - bich\_q(w\_m))/(1 - chi\_q(s\_m)), restated inside the proof at line 761, and Theorem cartan\_d (3.7) stating that c\textasciicircum{}q is a generalized Cartan matrix.
\end{itemize}\medskip
\ifdefined\hypertarget\hypertarget{sourcepointer-1}{}\fi
\noindent\textit{Source reference.} Original proposition 3.8: sources/sigma-2023-019-higher-braidings-of-diagonal-type/pinned-original-source.tex, source lines 734--734. Navigation pointer invoked at line 756: the counted text cites that numbered proposition of the same article, which is not retained in this card, so the reference is declared as a pointer carrying the number the article itself prints for it; the cited proposition is an allowed input and none of it is reproduced inside the span.\par\medskip
\ifdefined\hypertarget\hypertarget{sourcepointer-2}{}\fi
\noindent\textit{Source reference.} Original theorem 3.7: sources/sigma-2023-019-higher-braidings-of-diagonal-type/pinned-original-source.tex, source lines 714--714. Navigation pointer invoked at line 768: the counted text cites that numbered theorem of the same article, which is not retained in this card, so the reference is declared as a pointer carrying the number the article itself prints for it; the cited theorem is an allowed input and none of it is reproduced inside the span.\par\medskip
\begin{Proposition}\label{RdGraph}
Let $\dn\in \NN$ be even, $I=\{1,\dots,n\}$, and
$A$ be a set of tensors $\qq\in \KK^{I^\dn}$ with defined Cartan matrix $c^\qq$.
Assume that for all $\qq\in A$ and $\ell\in I$ we have
\[ \rho_\ell(\qq):=\sigma^{\qq}_\ell (\qq) \in A. \]
Then $(I,A,(\rho_i)_{i\in I},(c^\qq)_{\qq\in A})$ is a Cartan graph.
\end{Proposition}

\begin{proof}
Let $\qq\in A$ and $\ell,j \in I$.
The matrices $c^\qq$ are generalized Cartan matrices by Proposition~\ref{RdCartan}.
If Axiom (C2) is satisfied, then (C1) is satisfied as well because $\rho_\ell$ is defined using a reflection which has order two and $\sigma^{\qq}_\ell = \sigma^{\rho_\ell(\qq)}_\ell$ since the relevant Cartan entries are equal by~(C2).

For (C2), we need to compare $c^\qq_{\ell,j}$ and $c^{\rho_\ell(\qq)}_{\ell,j}$ for $\ell\ne j$.
With $m:=-c^\qq_{\ell,j}$, we have $R_m=0$ and $R_k\ne 0$ for $k<m$.
Now $R_m = (1-\bich_\qq(v_m))(1-\bich_\qq(w_m))/(1-\chi_\qq(s_m))$.
Recall that $v_m$ and $w_m$ are eigenvectors, i.e., $\big(\sigma^{\qq}_\ell\big)^{\otimes \dn}(v_m) = v_m$,
$\big(\sigma^{\qq}_\ell\big)^{\otimes \dn}(w_m) = -w_m$;
also, $\big(\sigma^{\qq}_\ell\big)^{\otimes \dn}(s_m) = s_m$ because $(m+1)s_m=v_m$.
Hence replacing $\qq$ by $\sigma^{\qq}_\ell (\qq)$ will map $R_m$ to
$(1-\bich_\qq(v_m))(1-\bich_\qq(-w_m))/(1-\chi_\qq(s_m))$.
But $1-\bich_\qq(-w_m)=0$ if and only if $-w_m=0$ or equivalently $w_m=0$.
Therefore, via~\eqref{cartan_d} the new condition produces the same Cartan entry
$c^\qq_{\ell,j}=c^{\rho_\ell(\qq)}_{\ell,j}$ for $\ell\ne j$.
\end{proof}

\end{document}
